\chapter{LITERATURE REVIEW}

\section{Digital Elevation Models (DEMs) for Terrain Representation}
Digital Elevation Models serve as the topographic foundation for terrain-based localization, representing surface relief through gridded elevation values \cite{purinton2017validation, okolie2024assessment}. Historically, the Shuttle Radar Topography Mission (SRTM) and ASTER GDEM provided the baseline for global coverage, but they often struggle with vertical inaccuracies in high-relief environments \cite{purinton2017validation}. Recent assessments indicate that Copernicus GLO-30 consistently outperforms these older models in mountainous regions, offering superior precision with a Root Mean Square Error (RMSE) of approximately 2.34m in high-relief sites \cite{okolie2024assessment}. While 10m resolution data is often viewed as a common baseline for urban or low-relief localization, research suggests that for the massive vertical scale of Himalayan peaks, 30m resolution is sufficiently representative \cite{purinton2017validation, okolie2024assessment}. This resolution offers a strategic balance, providing high-fidelity topographic fingerprints while significantly reducing the database size and computational overhead required for large-scale regional processing in Nepal.

\section{Horizon and Skyline-Based Geolocation}
Skyline matching identifies geographic coordinates by comparing silhouettes extracted from photographs with synthetic profiles rendered from a DEM \cite{baatz2012large, chen2015camera}. This approach is typically treated as a retrieval task against a pre-computed database of horizons, which is essential for navigation in environments where GPS is unreliable or ground-level imagery like Google Street View is sparse \cite{baatz2012large, rourera2022skyline}. A significant limitation in existing work is the sensitivity to ``bland'' or feature-poor skylines, which can lead to localization failure \cite{baatz2012large}. To combat this, the literature emphasizes the need for dense sampling of candidate viewpoints. The use of metric coordinate systems (e.g., UTM) for grid generation is recognized as a best practice to ensure uniform distance across rugged terrain, preventing the spatial distortions inherent in standard geographic coordinates when calculating line-of-sight in mountainous areas.

\section{Viewshed Analysis and Profile Generation}
Generating a reference database requires viewshed analysis to compute the visible horizon from a grid of candidate locations using ray-casting algorithms \cite{tzeng2013user, steger2022horayzon}. The computational cost of ray-casting is a function of the search radius and the angular resolution of the profile. Standard practices involve checking lines of sight up to 50km to capture distant high-altitude peaks that define the global horizon \cite{baatz2012large, steger2022horayzon}. While high-resolution profiles (e.g., 0.5 degrees) provide more detail, they increase the memory footprint of the searchable index. Efficient storage and retrieval are often managed using numerical fingerprints in optimized structures like NumPy arrays or FAISS (Facebook AI Similarity Search) indices, which allow for rapid querying across thousands of potential geographic candidates.

\section{Computer Vision for Horizon Extraction}
Automatic extraction of the skyline from raw images is a prerequisite for localization \cite{norinder2025panoramic, yang2025yunet}. Traditional methods utilizing Canny edge detection and gradient maps frequently fail under the haze, low-contrast, and cloudy conditions common in high-altitude photography \cite{baatz2012large, yang2025yunet}. Modern deep learning architectures offer increased robustness. Specifically, the MobileNet-V3-UNet architecture provides an efficient compromise between the high-pixel accuracy of traditional UNets and the low-latency requirements of mobile-ready systems \cite{yang2025yunet}. Despite the efficacy of deep learning, these networks require large datasets. In regions like the Himalayas where labeled ground truth is scarce, the literature supports a ``synthetic-to-real'' approach, where DEMs are used to render thousands of synthetic silhouettes for training, supplemented by simulated noise, clouds, and textures to improve model generalization.

\section{Similarity Matching and Orientation Estimation}
A major challenge in mountain geolocalization is the lack of accurate compass (heading) metadata. While astronomical cues like sun position (via ephemeris tools) can constrain the search space, they are weather-dependent \cite{baatz2012large, chen2015camera}. Consequently, the literature favors a rotational ``sliding window'' approach, comparing the extracted horizon against all possible orientations in the database \cite{chen2015camera}. For the matching process, a tiered strategy is often employed to maintain scalability: a fast initial search (Euclidean or Cosine similarity) identifies a subset of top candidates, which are then refined using Dynamic Time Warping (DTW) \cite{rourera2022skyline}. DTW is particularly valued for its ability to handle ``time-axis'' distortions caused by lens zoom or atmospheric warping. To keep the compute time linear, techniques such as Early Abandoning and Pruning (EAP) are critical for making DTW viable for real-time applications.

\section{Validation via Ground-Truth Datasets}
Validating terrain-based localization systems requires datasets with verified GPS and heading metadata. In the absence of massive hand-labeled datasets for the Nepal region, researchers frequently utilize Google Street View photospheres \cite{baatz2012large}. These provide a globally consistent ground-truth set that allows for the validation of both the horizon extraction accuracy and the similarity matching logic. This benchmarking is essential to ensure that the transition from synthetic training data to real-world photography does not result in a significant drop in localization precision.

\section{Conclusion and Project Gaps}
The existing research identifies several critical gaps: (i) the high computational cost of matching silhouettes without orientation data, (ii) the scarcity of labeled training data for the Himalayas, and (iii) the difficulty of maintaining accuracy in hazy conditions. This project addresses these gaps by implementing an efficient localization pipeline for Nepal that utilizes Copernicus GLO-30 DEM for topographic precision, a MobileNet-V3-UNet for robust haze-resistant segmentation, and a two-step similarity search (Euclidean + DTW) to ensure scalable, orientation-independent matching across the rugged Himalayan landscape.
